% GENERATED FILE. Edit canonical lessons, then run npm run generate:books.
% canonical-source-hash: fnv1a64:72ee332d53081678
% canonical-lessons: KA-C214-ancu, KA-C214-tudi, KA-C214-tala, KA-C214-addaraste, KA-C214-kavalu

\chapter{Edge, Tip, Bottom, Cross street, Fork}
\label{ch:ka-edge-tip-bottom-cross-street-fork}
\hlchaptermodality{214}

\begin{chapteropening}
\textbf{By the end of this chapter:} \emph{I can say an edge, a tip, the bottom, a cross street and a fork.}

Complete the last lesson of chapter 214: I can say an edge, a tip, the bottom, a cross street and a fork.
\end{chapteropening}

\section[\texorpdfstring{\kn{ಅಂಚು} (añcu)}{ಅಂಚು (añcu)}]{\kn{ಅಂಚು} (añcu) — an edge; the border of a sari}

\label{lesson:KA-C214-ancu}



\begin{quote}
Before the new one: say the Kannada for a stadium, then the Kannada for a museum.
\end{quote}

\subsection*{You'll want to know: \kn{ಅಂಚು}}
\textbf{\kn{ಅಂಚು}} — \emph{añcu} — \textquotedblleft{}an edge; the border of a sari\textquotedblright{}.

\begin{grammarlens}[title={What you've built}]
One more word for this chapter.
\end{grammarlens}

\subsection*{Your turn}
\begin{itemize}
\raggedright
  \item \emph{Say it:} \emph{añcu}
  \item \emph{Say it:} \emph{añcu}, once more
  \item \emph{Say it:} say \emph{añcu}
  \item \emph{Recall:} say the Kannada for a stadium, then the Kannada for a museum, then say \emph{añcu} again
\end{itemize}

\begin{tcolorbox}[breakable,colback=teal!4,colframe=teal!35!black,title={Before you move on}]
What does \kn{ಅಂಚು} mean? (\textquotedblleft{}An edge; the border of a sari\textquotedblright{}.) Say it once more.
\end{tcolorbox}
\section[\texorpdfstring{\kn{ತುದಿ} (tudi)}{ತುದಿ (tudi)}]{\kn{ತುದಿ} (tudi) — a tip, the far end}

\label{lesson:KA-C214-tudi}



\begin{quote}
Before the new one: say the Kannada for a museum, then the Kannada for an edge.
\end{quote}

\subsection*{You'll want to know: \kn{ತುದಿ}}
\textbf{\kn{ತುದಿ}} — \emph{tudi} — \textquotedblleft{}a tip, the far end\textquotedblright{}.

\begin{grammarlens}[title={What you've built}]
One more word for this chapter.
\end{grammarlens}

\subsection*{Your turn}
\begin{itemize}
\raggedright
  \item \emph{Say it:} \emph{tudi}
  \item \emph{Say it:} \emph{tudi}, once more
  \item \emph{Say it:} say \emph{tudi}
  \item \emph{Recall:} say the Kannada for a museum, then the Kannada for an edge, then say \emph{tudi} again
\end{itemize}

\begin{tcolorbox}[breakable,colback=teal!4,colframe=teal!35!black,title={Before you move on}]
What does \kn{ತುದಿ} mean? (\textquotedblleft{}A tip, the far end\textquotedblright{}.) Say it once more.
\end{tcolorbox}
\section[\texorpdfstring{\kn{ತಳ} (taḷa)}{ತಳ (taḷa)}]{\kn{ತಳ} (taḷa) — the bottom}

\label{lesson:KA-C214-tala}



\begin{quote}
Before the new one: say the Kannada for an edge, then the Kannada for a tip.
\end{quote}

\subsection*{You'll want to know: \kn{ತಳ}}
\textbf{\kn{ತಳ}} — \emph{taḷa} — \textquotedblleft{}the bottom\textquotedblright{}.

\begin{grammarlens}[title={What you've built}]
One more word for this chapter.
\end{grammarlens}

\subsection*{Your turn}
\begin{itemize}
\raggedright
  \item \emph{Say it:} \emph{taḷa}
  \item \emph{Say it:} \emph{taḷa}, once more
  \item \emph{Say it:} say \emph{taḷa}
  \item \emph{Recall:} say the Kannada for an edge, then the Kannada for a tip, then say \emph{taḷa} again
\end{itemize}

\begin{tcolorbox}[breakable,colback=teal!4,colframe=teal!35!black,title={Before you move on}]
What does \kn{ತಳ} mean? (\textquotedblleft{}The bottom\textquotedblright{}.) Say it once more.
\end{tcolorbox}
\section[\texorpdfstring{\kn{ಅಡ್ಡರಸ್ತೆ} (aḍḍaraste)}{ಅಡ್ಡರಸ್ತೆ (aḍḍaraste)}]{\kn{ಅಡ್ಡರಸ್ತೆ} (aḍḍaraste) — a cross street}

\label{lesson:KA-C214-addaraste}



\begin{quote}
Before the new one: say the Kannada for a tip, then the Kannada for the bottom.
\end{quote}

\subsection*{You'll want to know: \kn{ಅಡ್ಡರಸ್ತೆ}}
\textbf{\kn{ಅಡ್ಡರಸ್ತೆ}} — \emph{aḍḍaraste} — \textquotedblleft{}a cross street\textquotedblright{}.

\begin{grammarlens}[title={What you've built}]
One more word for this chapter.
\end{grammarlens}

\subsection*{Your turn}
\begin{itemize}
\raggedright
  \item \emph{Say it:} \emph{aḍḍaraste}
  \item \emph{Say it:} \emph{aḍḍaraste}, once more
  \item \emph{Say it:} say \emph{aḍḍaraste}
  \item \emph{Recall:} say the Kannada for a tip, then the Kannada for the bottom, then say \emph{aḍḍaraste} again
\end{itemize}

\begin{tcolorbox}[breakable,colback=teal!4,colframe=teal!35!black,title={Before you move on}]
What does \kn{ಅಡ್ಡರಸ್ತೆ} mean? (\textquotedblleft{}A cross street\textquotedblright{}.) Say it once more.
\end{tcolorbox}
\section[\texorpdfstring{\kn{ಕವಲು} (kavalu)}{ಕವಲು (kavalu)}]{\kn{ಕವಲು} (kavalu) — a fork (in a road)}

\label{lesson:KA-C214-kavalu}



\begin{quote}
Before the new one: say the Kannada for the bottom, then the Kannada for a cross street.
\end{quote}

\subsection*{You'll want to know: \kn{ಕವಲು}}
\textbf{\kn{ಕವಲು}} — \emph{kavalu} — \textquotedblleft{}a fork (in a road)\textquotedblright{}.

\begin{grammarlens}[title={What you've built}]
One more word for this chapter.
\end{grammarlens}

\subsection*{Your turn}
\begin{itemize}
\raggedright
  \item \emph{Say it:} \emph{kavalu}
  \item \emph{Say it:} \emph{kavalu}, once more
  \item \emph{Say it:} say \emph{kavalu}
  \item \emph{Recall:} say the Kannada for the bottom, then the Kannada for a cross street, then say \emph{kavalu} again
\end{itemize}

\begin{tcolorbox}[breakable,colback=teal!4,colframe=teal!35!black,title={Before you move on}]
What does \kn{ಕವಲು} mean? (\textquotedblleft{}A fork (in a road)\textquotedblright{}.) Say it once more.
\end{tcolorbox}
